\chapter{Homotopy limits and colimits: the practice}

\section{Convenient categories of spaces}

We now actually define various convenient categories of spaces that we have been referencing, such that we might obtain a complete cocomplete closed symmetric monoidal category, which is sometimes referred to as a ``cosmos'' since it is the best place to do enriched category theory. We go back to writing $\Top $ to denote the standard category of topological spaces. The main issue with this category is it's failiure to be closed under the cartesian monoidal structure:
\begin{theorem}\label{thm:6.1}
	The product, unit, and internal hom for any symmetric monoidal closed structure on $\Top $ necessarily have the cartesian product, singleton, and set of continuous functions as underlying sets. However, there is no suchc structure in which the monoidal product is the cartesian product, hence $\Top $ is not carteisna closed.
\end{theorem}

The details of the proof are that one cannot topologize the set $\Top (X,Y)$ such that evaluation $\Top (X,Y)\times X\to Y$ is continuous and satisfies the relevant universal property: $\Top (X,Y)$ should be terminal in the category of spaces $W$ with a continuous map $W\times X\to Y$. The situation is even worse for pointed spaces.

\begin{lemma}\label{lma:6.1}
	The smash product in $\Top _*$ is not associative. In particular, if $\mathbb{Q} $ and $\mathbb{N} $ inherit the subspace topology of $\mathbb{R} $ with basepoint 0, then $(\mathbb{Q} \wedge\mathbb{Q} )\wedge\mathbb{N} $ is not homeomorphic to $\mathbb{Q} \wedge (\mathbb{Q} \wedge\mathbb{N} )$.
\end{lemma}

We thus need to suggest alternatives to $\Top $, which contain as many of the important topological spaces as possible, in particular CW complexes, metric spaces, manifolds, and preferably even Hausdorff spaces.

\begin{definition}[Compactly Closed and $k$-Spaces]\label{dfn:6.1}
	A subspace $A\subseteq X$ is \emph{compactly closed} if for any continuous function $f:K\to X$ with $K$ compact Hausdorff, the preimage $f^{-1} (A)$ is closed in $K$. A space $X$ is called a $k$-space if every compactly closed subspace of $X$ is closed.
\end{definition}

We write $\mathrm{k} \Top $ for the full subcategory of $k$-spaces. The inclusion $\mathrm{k} \Top \hookrightarrow \Top $ admits a right adjoint $k$, called $k$-ification, which retopologizes the space $X$ such that closed subsets of $kX$ are precisely the compactly closed subsets of $X$. The identity $kX\to X$ is continuous and is universal in the sense that $kX$ is the finest topology such that $K\to X$ factors through $kX\to X$.

\begin{theorem}\label{thm:6.2}
	The category $\mathrm{k} \Top $ is ccomplete and cocomplete. Colimits are formed as in $\Top $, and limits are formed via applying $k$ to the corresponding limit in $\Top $. The category $\mathrm{k} \Top $ is cartesian closed.
\end{theorem}

The proof of cartesian closure is in an article I cannot access. The assertian about limits follows somehow idk. The main point is that the product of two $k$-spaces is the $k$-ification of the usual product topology. If one of these spaces is locally compact the topology is the same, but in general the $k$-ification is a finer topology. This modification means that geometric realization $\left| - \right| : \sSet \to \mathrm{k} \Top $ preserves products, which it didn't when the codomain was $\Top $.

\begin{lemma}\label{lma:6.2}
	If $X$ and $Y$ are simplicial sets, then the geometric realization of their product is the product of their geometric realization.
\end{lemma}

A function is $k$-continuous if and only if its composite with a continuous function $K\to X$ is continuous, for $K$ compact Hausdorff. Clearly a $k$-continuous function with domain as a $k$-space is continuous. The cartesian closed structure on $\mathrm{k} \Top $ is the $k$-ification of the compact-open topology, which we recall was defined as generated by
\[
	\left\{ \left\{ f : X \to Y\mid f(A)\subseteq U \right\} \mid A\in \mathcal{T}_X^{c} ,\qquad U\in \mathcal{T} _Y \right\} 
.\]
Here, $\mathcal{T} _X^{c} $ is the set of all compact subsets. The $k$-ification of this topology only differs if $X$ has non-Hausdorff compact subspaces.

$k$-continuity is nonlocal unless compact subsets of $X$ are locally compact, which is true if $X$ is Hausdorff. We have another assertion that fixes this issue, giving a better notion of convergence for sequences of continuous functions.

\begin{definition}[Weak Hausdorff/Compactly Generated]\label{dfn:6.2}
	A space $X$ is \emph{weak Hausdorff} if the image of any continuous map with compact Hausdorff domain is closed. A weakly Hausdorff $k$-space is called \emph{compactly generated}.
\end{definition}

Of course, all Hausdorff spaces are weak Hausdorff; we call Hausdorff $T_2$ and weak Hausdorff $T_1$. Examples of compactly generated spaces include metric spaces, CW complexes, and topological manifolds. Hence we obtain many of the important examples in topology. The idea is that the compact Hausdorff subspaces of a compactly generated space determine the topology. In particular, maps $X\to Y$ between compactly generated spaces are continuous if and only if the restriction to any compact Hausdorff subspaces of $X$ are continuous.

We write $\mathcal{C} \mathcal{G} \Top $ for the category of compactly generated spaces. The inclusion $\mathcal{C} \mathcal{G} \Top \hookrightarrow \mathrm{k} \Top $ has a left adjoint, which maps $K$-spaces to their maximal weak Hausdorff quotient. Important to note here is that this adjunction is monadic, so we obtain that $\mathcal{C} \mathcal{G} \Top \to \mathrm{k} \Top $ creates limits. Similarly, the left adjoint carries colimits down to $\mathcal{C} \mathcal{G} \Top $. Hence, limits formed in $\mathcal{C} \mathcal{G} \Top $ ar the same as in $\mathrm{k} \Top $, which is again applying $k$ to the ordinary limit. In particular, geometric realization is a strong monoidal functor $\left| - \right| : \sSet \to \mathcal{C} \mathcal{G} \Top $. Many colimits are weak Hausdorff already, but the ones that aren't are poorly behaved under the left adjoint $w$ of the function. Hence we only work in $\mathcal{C} \mathcal{G} \Top $ when the relevant colimits do not need weak Hausdorffication.

\begin{theorem}\label{thm:6.3}
	The categories $\mathrm{k} \Top $ and $\mathcal{C} \mathcal{G} \Top $ are cartesian closed. The categories $\mathrm{k} \Top_* $ and $\mathcal{C} \mathcal{G} \Top _*$ are closed symmetric monoidal with respect to the smash product.
\end{theorem}

\section{Simplicial model categories of spaces}

We now turn our attention to exploring simplicial homotopy structures and homotopy limits therein. The main example is that of $\sSet$. All objects are cofibrant, and fibrant objects are Kan complexes. These are simplicial sets in which \emph{all} horns admit fillers. We define the $i$th horn of the $n$-simplex $\Lambda_i^n$ as follows: the $j$th face map $d^j : \Delta^{n-1}  \to \Delta^{n} $ as a morphism of simplicial sets is defined as postcomposing with the coface map $\delta^j$. The image $\im(d^j)$ is a subsimplicial set of $\Delta^n$ defined levelwise as
\[
	\im(d^j)_m=\im(d^j_m)
.\]
Here we note that morphisms of simplicial sets have components for each $m$. The simplicial operators are simply the restrictions of those for $\Delta^n$ to the image. We then define the horn as the union
\[
	\Lambda_i^n \coloneqq \bigcup_{j\in [n],j\neq i} \im\left(d^j : \Delta^{n-1}  \to \Delta^n\right)
.\]
The union is once again defined levelwise, and since each simplicial set is a subsimplicial set of $\Delta^n$, we still have that the restriction of the simplicial operators is valid. The geometric intuition here is that the $j$th coface map inserts the standard $(n-1)$-simplex across from the $j$th vertex, hence we obtain the geometric notion of a horn via this construction. A Kan complex is then a simplicial set $X_{\bullet} $ in which all horns admit fillers, meaning for any diagram of solid arrows
\[\begin{tikzcd}[row sep = 2.5em, column sep = 2.5em]
\Lambda_i^n\ar[r] \ar[hook, d]  & X_{\bullet}  \\
\Delta^n\ar[ur, dotted]  & 
\end{tikzcd}\]
there exists a dotted arrow as indicated, rendering the diagram commutative. Note that all Kan complexes are weak Kan complexes, which require only that inner horns admit fillers (so $i\neq 0,n$), and hence are quasi-categories. Regular Kan complexes in fact define $\infty $-groupoids. I don't know why.

The usual geometric realization is strong monoidal as seen earlier in \cref{lma:6.2} when the codomain is restricted to a convenient category of spaces. Quillen showed that it is left Quillen, so we have that $\mathrm{k} \Top $ and $\mathcal{C} \mathcal{G} \Top $ are \emph{simplicial model categories}. We have proven part of this fact, see \cref{thm:3.1}. All objects are fibrant in these categories.

There is a strong monoidal Quillen adjunction with leftdjoint $(-)_+: \sSet \to \sSet_*$, as seen earlier. This makes based simplicial sets into a simplicial model category with all objects cofibrant, and fibrant objects as based Kan complexes. We also obtain

\begin{proposition}\label{prp:6.1}
	Geometric realization extends to a strong monoidal left Quillen functor $\left| - \right| : \sSet_* \to \mathcal{C} \mathcal{G} \Top _*$.
\end{proposition}

In particular, the monoidal categorical structures of $\mathcal{C} \mathcal{G} \Top_*$ and $\mathrm{k} \Top $ under the smash product are simplicial model structures. We can show that it is strong monoidal, however. The smash product is a pushout, and since geometric realization is a left adjoint, it preserves the pushout. It also commutes with the coproduct and product per above. Again, all objects are fibrant, and cofibrant objects are non-degenerately based cell complexes.

We have now lost all use for the category $\Top $, and for the purposes of this book (or at least this part), all limits, especially those used in the bar construction, are sufficiently well-behaved as to make the distinction between $\mathrm{k} \Top $ and $\mathcal{C} \mathcal{G} \Top $ immaterial. Hence all results proven are valid for either of them, so we write $\Top $ to denote either one of these categories.

\section{Warnings and simplifications}

\begin{remark}
	Given $\mathcal{M} $ homotopical, we always define the weak equivalences in $\mathcal{M} ^{\mathcal{D} } $ to be defined pointwise; hence weak equivalences are pointwise natural weak equivalences. I swap notation briefly to $\Fun$ instead of exponential to make the following diagram more clear. Universality of localization gives a map
	\[\begin{tikzcd}[row sep = 2.5em, column sep = 2.5em]
	\Fun(\mathcal{D} ,\mathcal{M} )\ar[" \Fun(\mathcal{D} \text{,} \gamma) ", r] \ar[" \gamma "', d]  & \Fun(\mathcal{D} ,\Ho \mathcal{M} ) \\
	\Ho (\Fun(\mathcal{D} ,\mathcal{M} ))\ar[dashed, ur]  & 
	\end{tikzcd}\]
	but this is only an equivalence of categories in very special cases, for example when $\mathcal{D} $ is discrete.

	The diagonal functor $\Delta : \mathcal{M}  \to \mathcal{M} ^{\mathcal{D} } $ is homotopical and is thus its own derived functor.\footnote{omg i finally figured out derived functors: its just the best approximation to make it homotopical! the total derived functor is the kan extension which is automatically homotopical but its between homotopy categories so like who cares, but the derived functor is a new functor that goes between the same categories as the original, but is homotopical, and agrees with the kan extension when postcomposed with delta and the kan extension is precomposed with gamma.} Since $\varinjlim \dashv\Delta\dashv \varprojlim  $, in particular the colimit, as seen before the adjunction promotes as
	\[
		\mathbf{L} \varinjlim \dashv \mathbf{L} \Delta \cong \Delta
	.\]
	Interestingly, this adjoint version of the diagonal functor may not be be the diagonal functor $\Ho \mathcal{M} \to \Ho (\mathcal{M} )^\mathcal{D} $, unless the comparison is an \emph{equivalence}. Hence, in absense of an equivalence, homotopy colimits need not be colimits in the homotopy category.
\end{remark}

Fix a field $k$. A chain complex $A_{\bullet} \in \Ch(k)$ is quasi-isomorphic to its homology, when differentials between the homologies are given as 0 morphisms. Among chain complexes with differentials given by 0, quasi-isomorphisms are just pointwise isomorphisms by definition of homology. Hence inverting quasi-isomorphisms in $\Ho (\Ch (k))$ is equivalent to the category of graded vector spaces, since a complex with 0 differentials is precisely a graded vector space. The latter is complete and cocomplete, but these (co)limits are not homotopy (co)limits. ``$\Ch(k)$ is a stable $(\infty ,1)$-category so homotopy pullbacks are also homotopy pushouts, which mutatis mutandis is not true for graded vector spaces.'' I don't know what this means, all the way from the remark about stable $\infty $-categories to ``mutatis mutandis''. Additionally, homotopy categories usually have very few limits and colimits. For example, not all coequalizers exist in $\Ho \Top $.

We provide the following simpliciation in most cases. By \cref{cor:5.1}, any $F : \mathcal{D}  \to \mathcal{M} $ with codomain a simplicial model category has $B(*,\mathcal{D} ,Q(F))$ has the right homotopy type for $\hocolim F$. However, if the diagram is already cofibrant, we have that $B(*,\mathcal{D} ,F)$ and $B(*,\mathcal{D} ,Q(F))$ are weakly equivalent by \cref{cor:5.2}, and the correction via $Q$ is not needed. The book rephrases this as: $Q$ simply corrects the homotopy type, but if it is already correct, we prefer a simpler construction.

\begin{remark}
	In $\Top $, viewed as a simplicial model cateogry, not all objects are cofibrant. They are all fibrant, so we can always compute the homotopy limit without the correction, but ``a folklore result'' shows that we can also compute homotopy colimits without correction as well. We will produce this proof later, after considering Reedy category theory more rigorously.
\end{remark}

\section{Sasmple homotopy colimits}
